\documentclass[10pt,a4paper]{amsart}
\usepackage[margin=19mm]{geometry}
\usepackage{amsmath,amssymb,mathtools}
\usepackage[scr=rsfs,cal=euler]{mathalfa}
\usepackage{xfrac}
\usepackage[unicode,hidelinks,bookmarks=false]{hyperref}
\newcommand{\ie}{i.e.\ }
\newcommand{\cf}{cf.\ }
\newcommand{\resp}{respectively\ }
\newcommand{\Subsection}{\subsection}
\providecommand{\nobreakdash}{\nobreak\hbox{-}}
\setlength{\emergencystretch}{2em}
\multlinegap0pt
\usepackage{bm}
\usepackage{bbm}
\usepackage{dsfont}
\usepackage{xspace}
\usepackage{cancel}
\usepackage{mathtools}
\usepackage{amsthm}
\makeatletter
\@ifundefined{claim}{\newtheorem{claim}{Claim}}{}
\@ifundefined{lemma}{\newtheorem{lemma}{Lemma}}{}
\@ifundefined{theorem}{\newtheorem{theorem}{Theorem}}{}
\makeatother
\newcommand\ACA{\mathsf{ACA}}
\newcommand\Nat{\mathbf{N}}
\newcommand\RAN{\mathsf{RAN}}
\newcommand\RCA{\mathsf{RCA}}
\newcommand\WOP{\mathsf{WOP}}
\newcommand\W{\mathrm{W}}
\newcommand\ap{\mathsf{ap}}
\newcommand\ordX{\mathcal{X}}
\newcommand\regHT{\mathsf{regHT}}
\title[Regressive versions of Hindman's Theorem]{Regressive versions of Hindman's Theorem}
\author{Lorenzo Carlucci, Leonardo Mainardi}
\date{Source: arXiv:2207.08554v4 (CC BY 4.0)}
\begin{document}
\maketitle

\noindent\emph{Source and licence.} Regressive versions of Hindman's Theorem. Version arXiv:2207.08554v4 (CC BY 4.0), distributed under the Creative Commons Attribution 4.0 International licence, as recorded in the arXiv licence metadata for this exact version and shown on the abstract page https://arxiv.org/abs/2207.08554v4. Full rights notice: licenses/arxiv-2207.08554-CC-BY-4.0.txt.\par\smallskip

\noindent\textit{Editorial scope.} Counted unit: the Theorem whose source label is \texttt{thm:reghtwop} in the article (source0.tex lines 1193--1196, source label \texttt{thm:reghtwop}), delivered together with the article's own complete original proof of it (source0.tex lines 1199--1314, one \texttt{proof} environment of 116 lines). No mathematics is written, repaired or re-derived: statement and proof are the article's own text line for line. Required context retained uncounted: thm:regtoaca (lines 971--975); claim1 (lines 1166--1169). Every definition, display and label of those lines is retained unchanged. Omitted from this package and not redistributed: 1--970, 976--1165, 1170--1192, 1197--1198, 1315--1565. Nothing inside a retained excerpt is omitted or altered: every retained body is delivered exactly as the article's lines are written, so no omission receipt of any kind accompanies this package. Citations: the retained excerpts carry no citation command at all, so no bibliography entry of the article is transcribed and no reference is redistributed. Reference adaptation: none was needed and none was made; every reference inside the delivered unit resolves inside it, so no unresolved reference or citation is produced. Printed slips kept literal: no printed word, symbol, hypothesis or number of the counted statement or of its proof was altered. Layout and class: the article's own class is \texttt{article} with packages including authblk, inputenc, babel, fullpage, xspace, amsmath, amsthm, amssymb, dsfont, tikz-cd, enumitem, hyperref; the wrapper is the lane's pinned \texttt{amsart} editorial wrapper at 10pt on A4, which restates the theorem-like environments the retained text needs and adds the article's own macro definitions that the retained text needs (\texttt{\textbackslash ACA}, \texttt{\textbackslash Nat}, \texttt{\textbackslash RAN}, \texttt{\textbackslash RCA}, \texttt{\textbackslash W}, \texttt{\textbackslash WOP}, \texttt{\textbackslash ap}, \texttt{\textbackslash ordX}, \texttt{\textbackslash regHT}), with extra wrapper packages (bm, bbm, dsfont, xspace, cancel, mathtools). The delivered numbering is the unit's own. Assets: no retained span contains a figure environment, a graphics-inclusion command, a PDF-page inclusion, a table or a TikZ picture; no image file is copied into this package, the package declares no asset dependency and no image file is redistributed with it. Licence: version arXiv:2207.08554v4 of this article is distributed under the Creative Commons Attribution 4.0 International licence, as recorded in the arXiv licence metadata for this exact version and shown on the abstract page https://arxiv.org/abs/2207.08554v4. A full rights notice is delivered at licenses/arxiv-2207.08554-CC-BY-4.0.txt. Characters: every retained excerpt is pure ASCII, so the delivered engine, XeLaTeX with its default Latin Modern text font, reports no missing character.
\begin{theorem}\label{thm:regtoaca}
Let $n\geq 2$.
$\lambda\regHT^{= n}[\ap]$ implies $\ACA_0$ over $\RCA_0$. Moreover, 
$$\lambda\regHT^{=n}[\ap]\geq_\W \RAN.$$
\end{theorem}

\begin{lemma}\label{claim1} The following is provable in $\RCA_0$: If $\alpha=(\alpha_i)_{i\in\Nat^+}$ is an infinite descending sequence in $\omega^\ordX$, then
$$\forall n\ \exists n'\ \exists m\leq |\alpha_n| \ \big(n'>n \ \land m \leq |\alpha_{n'}| \ \land \alpha_{n,m} >_{\ordX} \alpha_{n',m}\big),$$
where $\alpha_{i,j}$ denotes the $j$-th component of $\alpha_i$ for $j \in [1,|\alpha_i|]$ and is otherwise undefined.
\end{lemma}

\begin{theorem}\label{thm:reghtwop}
Let $n \geq 2$. $\lambda\regHT^{= n}[\ap]$ implies $\WOP(\ordX\mapsto \omega^\ordX)$ over $\RCA_0$.
Moreover, $\lambda\regHT^{=n}[\ap]\geq_\W \WOP(\ordX\mapsto \omega^\ordX)$.
\end{theorem}

\begin{proof}
Let $\alpha= (\alpha_n)_{n\in \Nat^+}$ be an infinite descending sequence in $\omega^\ordX$. 
%We denote by $\alpha_{n,m}$ the $m$-th exponent of $\alpha_n$, where $1 \leq m \leq |\alpha_n|$.
We say that $\alpha_{n,m}$ is \emph{decreasible} if there exists a $n'>n$ such that $\alpha_{n',m} <_\ordX \alpha_{n,m}$. In this case we say that $\alpha_{n',m}$ \emph{decreases} $\alpha_{n,m}$. 
With this terminology Lemma \ref{claim1} says that $\RCA_0$ knows that for all $i\geq 1$ there exists $j \in [1, |\alpha_i|]$ such that $\alpha_{i,j}$ is decreasible. If $\alpha_{n',m}$ decreases $\alpha_{n,m}$ and no $\alpha_{k,m}$ 
with $k < n'$ decreases $\alpha_{n,m}$ we call $\alpha_{n',m}$ the \emph{least decreaser} of $\alpha_{n,m}$. 

Now suppose that $f:\Nat \to \Nat$ is a function with the following property:

\bigskip
{\em Property P}: For all $i\in \Nat^+$ for all $j \in [1, |\alpha_i|]$
if $\alpha_{i,j}$ is decreasible then $\alpha_{i,j}$ is decreased by $\alpha_{k,j}$ for some $k\leq f(i)$. 

\bigskip
We first show that given such an $f$ we can compute (in $f$ and $\alpha$) 
an infinite descending sequence $(\sigma_i)_{i\in\Nat^+}$ in $\mathcal{X}$ as follows. 

\bigskip
{\bf Step $1$}.  
Pick the leftmost decreasible component of $\alpha_1$ (which exists by Lemma \ref{claim1}). 
This can be done by inspecting all components 
in $\alpha$ up through $\alpha_{f(1)}$, since $f$ has Property $P$. 

Let $p_1$ be the position of the leftmost decreasible component of $\alpha_1$. 
Pick the smallest $d_1 \leq f(1)$ such that $\alpha_{d_1, p_1}$ decreases 
$\alpha_{1,p_1}$. 
We set $\sigma_1 = \alpha_{d_1, p_1}$ and observe that
all decreasible components in $\alpha_{d_1}$ occur at positions $\geq p_1$. 
Suppose otherwise and let $1\leq p^* < p_1$ be such that $\alpha_{d_1, p^*}$ is decreasible. 
Let $d^* > d_1$ such that $\alpha_{d^*, p^*}$ decreases $\alpha_{d_1,p^*}$. Then 
$\alpha_{d^*, p^*} <_{\ordX} \alpha_{d_1,p^*}$ by definition of decreasible. 
On the other hand, by choice of $d_1$ and $p_1$, and since
$p^* < p_1$, it must be the case that $\alpha_{d_1, p^*} = \alpha_{1, p^*}$. Hence $\alpha_{1, p^*}$ is a decreasible
component in $\alpha_1$ on the left of position $p_1$, which contradicts the choice of $p_1$. 

\bigskip
{\bf Step $i+1$ ($i > 0$)}. Suppose $d_i, p_i, \sigma_i$ are defined so that $\sigma_i = \alpha_{d_i, p_i}$, $(\sigma_j)_{1\leq j\leq i}$ is decreasing in $\mathcal{X}$ and all decreasible components in $\alpha_{d_i}$ occur at positions $\geq p_i$. 

Pick the leftmost decreasible component in $\alpha_{d_i}$ (which exists by Lemma \ref{claim1}). This can be done by inspecting all components in $\alpha$ up to 
$\alpha_{f(d_i)}$, since $f$ has Property $P$.  Let $\alpha_{d_i, \ell}$ be the chosen component. Set $p_{i+1}=\ell$ and note that necessarily $p_{i+1}\geq p_i$.

%Notice that $\alpha_{d_i}$ contains no decreasible term in position $< p_{i+1}$. Since $\alpha$ is decreasing this implies, in particular, that $\alpha_{d_i, k} = \alpha_{d^*,k}$ for all $d^* > d_i$ and all $k<p_{i+1}$.

Pick $d \leq f(d_i)$ minimal such that $\alpha_{d, p_{i+1}}$
decreases $\alpha_{d_i, p_{i+1}}$. Set $d_{i+1}=d$. 
Let $\sigma_{i+1} = \alpha_{d_{i+1}, p_{i+1}}$. 
Obviously $\sigma_i >_\mathcal{X} \sigma_{i+1}$, since $\sigma_i  = \alpha_{d_i,p_i} \geq \alpha_{d_i,p_{i+1}} >_\mathcal{X} \alpha_{d_{i+1}, p_{i+1}} = \sigma_{i+1}$ (note that $p_i \leq p_{i+1}$).

We observe that also the last part of the inductive invariant is guaranteed, since no decreasible component 
in $\alpha_{d_{i+1}}$ occurs on the left of $p_{i+1}$. 
Suppose otherwise as witnessed by $1\leq p^* < p_{i+1}$. Let $d^* > d_{i+1}$ such that $\alpha_{d^*,p^*}$
decreases $\alpha_{d_{i+1}, p^*}$. Then $\alpha_{d^*, p^*}$ also decreases $\alpha_{d_i, p^*}$ 
since $\alpha_{d_{i+1}, p^*}= \alpha_{d_{i}, p^*}$, where the latter is due to the fact that $\alpha$ is decreasing 
and $p^*$ is less than $p_{i+1}$, which is the position of the leftmost decreasible component in $\alpha_{d_i}$. This contradicts the choice of $p_{i+1}$. 

\bigskip
We next show how to obtain a function satisfying Property $P$ from a solution of $\lambda\regHT^{=n}[\ap]$ for a suitable colouring. The argument is similar to the proof of Theorem \ref{thm:regtoaca}.

For this purpose it is convenient to use a sequence $\beta$ of all the components of the terms $\alpha_n$ in $\alpha$, enumerated in order of appearance: more precisely, $(\beta_h)_{h\in\Nat^+}$ is the ordered sequence $\alpha_{1,1}, \alpha_{1,2}, \dots, \alpha_{1,|\alpha_1|}, \alpha_{2,1}, \alpha_{2,2}, \dots, \alpha_{2,|\alpha_2|}, \dots$. 
This sequence is obviously easily computable from $\alpha$. 
Formally we construct such a sequence by first defining a function $\iota : \Nat^+\times\Nat^+\to \Nat^+$ as follows: $\iota(n,m)= \sum_{1\leq k < n} |\alpha_k| + m$, for all $n\in\Nat^+$ and all $m \in [1, |\alpha_n|]$, while 
$\iota(n,m)=0$ in all other cases. 
%The function $\iota$ is bijective on the set $\{ (n,m) \,:\, n\in\Nat^+ \mbox{ and } 1 \leq m \leq |\alpha_n| \}$. 
We correspondigly fix functions $t: \Nat^+ \to \Nat^+$ and $p:\Nat^+\to\Nat^+$ such that for each $n\in\Nat^+$ we have
$\iota(t(n),p(n))=n$. The sequence $(\beta_h)_{h\in\Nat^+}$ of all components appearing in $\alpha$ is then defined
by setting $\beta_h = \alpha_{t(h),p(h)}$. 


Define $c:\Nat\to \Nat$ as follows: $c(x)=$ the unique $i < \lambda(x)$ satisfying the following conditions:
\begin{enumerate}
\item There exists $j$ such that $\lambda(x) \leq j < \mu(x)$ and $\beta_j$ is the least decreaser of $\beta_i$, and 
\item For all $j'$ such that $j < j'  < \mu(x)$, if $\beta_{j'}$ is the least decreaser of $\beta_{i'}$ then $i' \geq \lambda(x)$. 
\end{enumerate}
If no such $i$ exists, we set $c(x)=0$.

The function $c$ is computable in $\alpha$ and $\lambda$-regressive. Let $H = \{ h_1 < h_2 < h_3 < \dots \}$ be an apart solution to $\lambda\regHT^{= n}$ for $c$. The following Claim ensures the existence of an $(\alpha \oplus H)$-computable function with Property $P$.

\begin{claim}\label{claim2} For each $h_k\in H$ and each $\alpha_{\ell,m}$ such that $\iota(\ell,m) < \lambda(h_k)$, if there exists $\alpha_{\ell',m}$ such that $\alpha_{\ell',m}$ decreases $\alpha_{\ell,m}$ then there exists such an $\alpha_{\ell',m}$ with $\iota(\ell',m)<\mu(h_{k+n-1})$.
\end{claim}

\begin{proof}[Proof of Claim \ref{claim2}]
Assume by way of contradiction that there is some $h_k \in H$ and some $\alpha_{\ell,m}$ with $\iota(\ell,m)=i < \lambda(h_k)$ such that $\alpha_{\ell,m}$ is decreasible but not by any $\alpha_{\ell',m}$ with $\iota(\ell',m) < \mu(h_{k+n-1})$. 

Let $b$ be such that if $\alpha_{\ell'',m}$ is decreasible and $\iota(\ell'',m) < \lambda(h_k)$, then there exists $\ell',m$ such that $\iota(\ell',m) < b$ and $\alpha_{\ell',m}$ decreases $\alpha_{\ell'',m}$. The existence of $b$ can be proved in $\RCA_0$ using the following instance of strong $\Sigma^0_1$-bounding (similarly as in the proof of Theorem \ref{thm:regtoaca}):
$$ \forall n \exists b \forall i < n ( \exists j (\alpha_{t(j),p(j)} \mbox{ decreases } \alpha_{t(i),p(i)}) \to \exists j < b (\alpha_{t(j),p(j)} \mbox{ decreases } \alpha_{t(i),p(i)}).$$
Since $H$ is infinite, there is an $h_{k'} \in H$ such that $h_{k'} > h_{k+n-1}$ and $\mu(h_{k'}) \geq b$. Then, by min-term-homogeneity, $c(h_k + \cdots + h_{k+n-1}) = c(h_k + \cdots + h_{k+n-2} + h_{k'})$. But by choice of $h_k$, $h_{k'}$ and the definition of $c$, we can show that $c(h_k + \cdots + h_{k+n-1}) \neq c(h_k + \cdots + h_{k+n-2} + h_{k'})$, 
yielding a contradiction.

To see this we reason as follows. First observe that, by apartness of $H$, the following identities hold:
$$\lambda(h_k + \cdots + h_{k+n-1}) = \lambda(h_k + \cdots + h_{k+n-2} + h_{k'}) = \lambda(h_k),$$
and
$$\mu(h_k + \cdots + h_{k+n-2} + h_{k'})=\mu(h_{k'}).$$
Let $j \in [\mu(h_{k+n-1}), \mu(h_{k'}))$ be such that $\alpha_{t(j),p(j)}$ is the least decreaser of $\alpha_{\ell,m}$. Such a $j$ exists by choice of $\alpha_{\ell,m}$, $h_k$ and $h_{k'}$. In fact, by hypothesis, 
$\alpha_{\ell,m}$ is decreasible but not by any component with $\iota$-index below $\mu(h_{k+n-1})$. By choice
of $h_k'$ the least decreaser of $\alpha_{\ell,m}$ must have $\iota$-index smaller than $\mu(h_{k'})$, since $\iota(\ell,m) < \lambda(h_k)$.

First note that
$c(h_k + \cdots + h_{k+n-2} + h_{k'})$ cannot be $0$, 
since this occurs if and only if there is no $i^* < \lambda(h_k)$ such that for some $j^* \in [\lambda(h_k), \mu(h_{k'}))$, $\alpha_{t(j),p(j)}$ decreases $\alpha_{t(i^*), p(i^*)}$; but the latter is false by choice of $h_k$ and $h_{k'}$. 

If $c(h_k + \cdots + h_{k+n-1})$ takes some non-zero value $i^* < \lambda(h_k)$, then this same value cannot
be taken by $c(h_k + \cdots + h_{k+n-2} + h_{k'})$ under our assumptions. 
If it were, it would mean that $\alpha_{t(i^*),p(i^*)}$ is decreased for the 
first time by some $\alpha_{t(j^*), p(j^*)}$ with $j^* < \mu(h_{k'})$ such that $j^*$ is also maximal below $\mu(h_{k'})$ such that $\alpha_{t(j^*), p(j^*)}$ is the least decreaser of some $\alpha_{t(q),p(q)}$ with  $q < \lambda(h_k)$. This is impossible since the least decreaser of $\alpha_{t(i^*),p(i^*)}$, by the hypothesis that $c(h_k + \cdots + h_{k+n-1})=i^*$, occurs earlier in the sequence of the $\beta_h$'s than the least decreaser of $\alpha_{t(i),p(i)}$ since, by the definition of $c$, it must be that $j^* < \mu(h_k + \cdots + h_{k+n-1})$ and the latter value, by apartness, equals $\mu(h_{k+n-1})$, as noted above. On the other hand, $j$ is in $[\mu(h_{k+n-1}), \mu(h_{k'}))$, so that $j^*<j$. Thus $j^*$ cannot be maximal below $\mu(h_{k'})$ such that $\alpha_{t(j^*),p(j^*)}$ is the least decreaser of some $\alpha_{\ell'',m}$ with $\iota(\ell'',m)$ below $\lambda(h_k)$, as required by the definition of $c$, since $\alpha_{t(j),p(j)}$ is such a least decreaser of $\alpha_{t(i),p(i)}$, and $i<\lambda(h_k)$. 

This proves the Claim.
\end{proof}

% Conclusion of the argument
Now it is sufficient to observe that the $(\alpha \oplus H)$-computable function 
$f$ defined as follows has the Property $P$: on input $n$, 
pick the least $k$ such that $\sum_{1\leq n' \leq n} |\alpha_{n'}| < \lambda(h_k)$ and 
let $f(n)$ be the $\alpha$-index of the $\mu(h_{k+n-1})$-th 
element in the sequence $\beta$ of all components appearing in $\alpha$, i.e., $f(n) = t(\mu(h_{k+n-1}))$. 
That this choice of $f$ satisfies Property P is implied by Claim \ref{claim2} above. This concludes the proof of the theorem.
\end{proof}

\end{document}
